\Needspace{6\baselineskip}
\begin{center}\hypertarget{native-262-article189}{}\pdfbookmark[2]{189.}{article-189}189.\end{center}
Dieses Verfahren unterscheidet sich vom vorangegangenen nur darin, dass
es zwei geozentrischen Örtern genau und den übrigen möglichst genau
entspricht. Beim anderen Verfahren wird keine Beobachtung gegenüber den
übrigen bevorzugt, und die Fehler werden so weit wie möglich auf alle
verteilt\rlap{\raisebox{1.4ex}{\rule{.6em}{.35pt}}}. Das Verfahren des vorigen
Artikels steht daher nur insofern hinter dem früheren zurück, als sich
die Fehler bei den übrigen Örtern merklich verringern lassen, wenn man
einen Teil der Fehler den Örtern $L,L'$ zuweist. Meist lässt sich jedoch
durch eine geeignete Wahl der Beobachtungen $L,L'$ leicht verhindern,
dass dieser Unterschied ins Gewicht fällt. Man muss also darauf
\SourceLeaf{263}achten, für $L,L'$ Beobachtungen von besonders hoher Genauigkeit
zu wählen, die außerdem so beschaffen sind, dass kleine Änderungen der
geozentrischen Positionen selbst die aus ihnen und den Entfernungen
abgeleiteten Elemente nicht allzu stark beeinflussen. Es wäre daher
unzweckmäßig, Beobachtungen zu wählen, die zeitlich dicht beieinanderliegen
oder denen nahezu entgegengesetzte oder nahezu zusammenfallende
heliozentrische Örter entsprechen.
\begin{center}\rule{3cm}{.4pt}\end{center}
